% Transcription — fonds Alexandre Grothendieck, shelfmark 161-2, batch 4 (pages 61-80).
% Established from the facsimile archives/batches/161-2/batch-04.pdf (a mirror of
% https://grothendieck.umontpellier.fr/161-2.pdf), per the skill
% transcribe-grothendieck. Page 79 is blank and absent; page 66 carries a single
% opening bracket, kept. The batch is one run of yellowed and white lined sheets
% on algebraic theories over a category Δ of categories with prescribed limits:
% p. 61 closes the proof on accessible objects begun in batch 3, then
% « Compléments » (pp. 62-68: transfinite construction of theories, exactness
% conditions after Schanuel, change of Δ, limits in T(C), the embedding
% R° ↪ S = T(Ens)), « Exemples » (pp. 69-78: the cases d' = d'' = ∅, no
% exactness conditions, finite limits, presentation-finite objects and
% Ind(Σ) ≃ S keyed to SGA 4 I 8.7-8.9), and a closing list of the cases to treat
% (p. 80).
% Pass: Fable 5.1 (claude-fable-5-1), 2026-09-04 — first pass, unchecked against the pages by a human.
\documentclass[11pt,a4paper]{article}
% Le préambule commun à toutes les transcriptions du fonds Grothendieck.
%
% Il tient en une idée : **l'incertitude se compose, elle ne se résout pas.**
% Une transcription qui lisse un mot douteux a détruit la seule chose pour
% laquelle on la faisait — un lecteur doit pouvoir savoir, à chaque endroit,
% ce qui était sur le feuillet et ce qui a été ajouté. Les six macros qui
% suivent sont donc l'appareil critique, et non de la décoration.
%
% Les mêmes commandes sont reconnues par `scripts/render.mjs`, qui produit la
% vue de lecture du site. Une macro ajoutée ici doit l'être aussi là-bas,
% faute de quoi le rendu échouera bruyamment — ce qui est voulu : un rendu
% silencieusement faux est pire qu'un rendu absent.

\ProvidesPackage{grothendieck}[2026/08/08 Transcriptions du fonds Grothendieck]

\usepackage{amsmath,amssymb,amsthm}
\usepackage{mathrsfs}
\usepackage{stmaryrd}
% Les diagrammes commutatifs sont partout dans ce fonds : c'est la langue
% même de la géométrie algébrique de Grothendieck, et une transcription qui
% les rendrait en prose ne transcrirait rien.
\usepackage{tikz-cd}
% Français par défaut — c'est la langue des feuillets — mais l'anglais est
% chargé aussi : la traduction et le résumé sont des documents anglais, et la
% typographie française (espaces avant les deux-points) y serait une faute.
% Ces documents basculent par \selectlanguage{english} en tête de corps.
\usepackage[english,french]{babel}
\usepackage{fontspec}
\usepackage{geometry}
\geometry{a4paper,margin=25mm,marginparwidth=28mm}
\usepackage{marginnote}
\usepackage{xcolor}
% ulem, not soul. `soul`'s \st and \ul are built on a character-by-character
% scan that XeLaTeX's font machinery breaks: the document fails with an
% undefined control sequence rather than degrading. ulem does the same two jobs
% and is Unicode-engine safe. `normalem` stops it hijacking \emph, which the
% transcriptions use for Grothendieck's own emphasis.
\usepackage[normalem]{ulem}
\usepackage{hyperref}
\hypersetup{colorlinks=true,linkcolor=black,urlcolor=black,pdfborder={0 0 0}}

% --- Métadonnées du lot -----------------------------------------------------
% Reprises telles quelles par le script de rendu, qui les affiche en tête de
% la vue de lecture. Elles fixent ce que le fichier prétend couvrir : sans
% elles, une transcription flotte sans point d'attache dans un fonds de seize
% mille feuillets.

\newcommand{\folder}[1]{\gdef\grothFolder{#1}}
\newcommand{\batch}[1]{\gdef\grothBatch{#1}}
\newcommand{\leaves}[2]{\gdef\grothFirst{#1}\gdef\grothLast{#2}}
\newcommand{\foldertitle}[1]{\gdef\grothFoldertitle{#1}}
\gdef\grothFolder{?}\gdef\grothBatch{?}\gdef\grothFirst{?}\gdef\grothLast{?}\gdef\grothFoldertitle{}

% --- Le résumé --------------------------------------------------------------
%
% Il ouvre la lecture modernisée, sous le titre « Résumé », et il est composé
% en petit. Ce n'est pas un ornement : c'est le seul passage du document qui
% ne soit pas des mathématiques, et un lecteur doit pouvoir le reconnaître
% avant de l'avoir lu — puis le sauter s'il connaît déjà le sujet, ou s'y
% arrêter s'il ne le connaît pas. Le filet qui le suit marque la reprise du
% corps.

\newenvironment{resume}
  {\small\setlength{\parskip}{0.45em}}
  {\par\medskip\hrule\medskip}

% --- Mots-clés ---------------------------------------------------------------
%
% En anglais, en clôture du résumé de la lecture modernisée : le vocabulaire
% moderne sous lequel chercher ce que les feuillets construisent. Le manifeste
% du site (`npm run manifest`) les extrait pour étiqueter la cote entière —
% cette ligne est la source unique des tags, il n'y a pas de fichier à côté.
\newcommand{\keywords}[1]{\par\smallskip\noindent
  {\footnotesize\textsc{Keywords} --- \textit{#1}}\par}

% --- L'appareil critique ----------------------------------------------------

% Le numéro de feuillet, en marge. C'est l'ancre qui permet de régler un
% désaccord sur une formule en regardant *un* feuillet plutôt qu'une chemise
% de six cents. Le site s'en sert aussi pour tourner les pages du fac-similé
% quand on fait défiler la transcription.
\newcommand{\leaf}[1]{%
  \marginnote{\footnotesize\color{gray!70}\textsf{#1}}%
  \label{leaf:#1}\ignorespaces}

% Illisible : on ne devine pas. Un mot inventé qui se lit comme les autres est
% le pire résultat possible.
\newcommand{\ill}{\textcolor{red!60!black}{[\,\ldots\,]}}

% Lecture douteuse : le mot est donné, et signalé comme incertain.
\newcommand{\uncertain}[1]{\uline{#1}}

% Ajout de l'éditeur — une lettre manquante, un mot sous-entendu.
\newcommand{\add}[1]{\textcolor{blue!50!black}{[#1]}}

% Biffé par Grothendieck lui-même. Ce qu'il a rayé fait partie du document :
% une rature dit souvent où la pensée a changé de direction.
\newcommand{\struck}[1]{\sout{#1}}

% Note du transcripteur, sur l'état matériel du feuillet.
\newcommand{\note}[1]{\par\smallskip{\small\color{gray!60!black}\textit{[#1]}}\par\smallskip}

% Note marginale de Grothendieck, distincte du corps du texte.
\newcommand{\marginal}[1]{\par\smallskip\noindent\hspace{1em}%
  {\small\color{gray!60!black}#1}\par\smallskip}

% --- Titre ------------------------------------------------------------------

\AtBeginDocument{%
  \begin{center}
    {\large\bfseries Fonds Alexandre Grothendieck --- cote n\textsuperscript{o}~\grothFolder}\\[2pt]
    {\itshape\grothFoldertitle}\\[4pt]
    {\small feuillets \grothFirst--\grothLast\ (lot \grothBatch)}\\[2pt]
    {\footnotesize\color{gray!60!black}Transcription établie d'après le fac-similé numérisé
     par l'Université de Montpellier.\\
     Les lectures incertaines sont soulignées, les illisibles notées [\,\ldots\,],
     les ajouts entre crochets.}
  \end{center}
  \vspace{1em}}

\endinput


\folder{161-2}
\batch{4}
\pages{61}{80}
\dating{[vers 1963-1973]}
\watermark{Édition de démonstration}
\foldertitle{Algèbre universelle [ou catégories] : notes manuscrites (s.d.)}

\begin{document}

\page{61}
\note{la démonstration du corollaire qui ferme la p. 60 ; les $Y$, $C$,
$E_{\pi}$ sont ceux de la p. 60}
\underline{Dém}\ Soit $\pi$ un cardinal \add{infini} tel que les objets de
$C$ soient $\pi$-accessibles. Je veux à prouver que la sous-catégorie
$E_{\pi}$ de $E$ formée des objets $\pi$-accessibles est \struck{\ill{}}
équiv. à une catégorie petite. \struck{Je sais} Sps $C$ stable par limites
inductives finies (sinon on \uncertain{s'y ramène}, $C$ reste
\struck{\ill{}} petite). Soit $X$ un objet $\pi$-accessible, on a
$X = \varinjlim_{C/X} Y$, et $C/X$ est filtrant. \add{(on a donc aussi
$X = \varinjlim_{I} X_{i}$, où $I$ \ill{} \ill{} ordonné filtrant, les
$X_{i} \in C$)} Soit $J$ l'un des \struck{sous-catégories} \add{ensembles
ordonnés} filtr. \ill{} de $C/X$ de cardinal $\leqslant \pi$, alors $J$
est \struck{\ill{}} grand devant $\pi$, et on a
\[
  X = \varinjlim_{j \in J} Y_{j}, \qquad
  Y_{j} = \varinjlim_{I_{j}} \struck{\ill{}} Y_{i} .
\]
\note{la seconde formule est écrite au-dessous de la première ; la lecture
des indices est douteuse}
Comme $X$ est $\pi$-accessible, \struck{\ill{}} l'identité
$X \xrightarrow{\ \sim\ } \varinjlim_{j \in J} Y_{j} = X$ se factorise par
un $Y_{j}$, donc $X$ est isom à un facteur direct d'un $Y_{j}$. Donc tout
objet $\pi$-accessible est isom à un facteur direct d'un objet
$\varinjlim_{i \in I} Y_{i}$ (où $I$ \uncertain{ens. ordonné}) de card
$\leqslant \pi$, \add{relevant les $Y_{i}$ dans $C$}. \struck{Ceci donne
le résultat \ill{}} \note{une ligne et demie biffées, avec des mots
surajoutés qui ne se laissent pas lire}

\note{bloc encadré, relié par un trait au passage précédent :}
\struck{$Y_{i,n} = Y_{i}$, $Y_{i,m}$, $Y_{j,m}$} trouver \struck{que} on
$\simeq$ \ill{} pour tt objet de $C$ (Cat stable par $\varinjlim$
finies) $\pi$-accessible, \add{limites des} \struck{systèmes inductifs
filtrants} $E_{\pi}$ est l'ensemble des foncteurs $\varinjlim_{I} Y_{i}$,
les $Y_{i} \in C$, card $I \leqslant \pi$, \struck{\ill{}}

On va en conclure que \struck{tt objet} $X$ est \struck{\ill{}} limi\ill{}
\ill{} une telle limite, ou encore que la catégorie filtrante $C/X$ admet
un ensemble cofinal de card $\leqslant \pi$. \struck{Donc \ill{} \ill{},
la $Y$ \ill{} \ill{}} Je dis que \struck{\ill{}} sous-catégorie
$F^{\pi}_{\uncertain{\pi}}$ de $E$ formée des $\varinjlim_{i \in I} Y_{i}$,
avec $Y_{i} \in C$, $I$ ens. ordonné tel que Card $I \leqslant \pi$,
est stable par $\varinjlim$ de cardinal $\leqslant \pi$. Il suffit de
prouver que c'est stable par facteurs directs (car \ill{} \ill{}) puis
aussi alors que $F^{\pi} = E_{\pi}$, et \ill{} \uncertain{-pliques} SGA 4
I 9.9. Donc il suffit de prouver que c'est stable par $\varinjlim$
filtrantes dénombrables (disons) \ldots

\section{Compléments}
\note{titre souligné en tête de la p. 62, de sa main ; les alinéas sont
numérotés de 1 à 4 dans des cercles, en marge}

\page{62}
\begin{enumerate}
  \item[(1)] \underline{Pb}\ Pour $R \in \mathrm{Ob}\,\Delta$ et
    $I \to R$, \add{$I$ cat. discr.} caractériser la situation où cela
    correspond à une famille d'objets de base (indexée par $I$) pour une
    théorie \underline{algébrique} $T_{R}$. (Il y a une caractérisation
    constructive transfinie : $R$ est l'aboutissement d'une suite
    transfinie de catégories $R_{\alpha}$ sur $R$, faisant une (pseudo)
    \struck{\ill{}} objet inductif de catégories \struck{\ill{} dans
    $\Delta$, avec}
    \begin{enumerate}
      \item[(1)] $R_{0} = I$
      \item[(2)] $R_{\alpha+1}$ déduit de $R_{\alpha}$ soit en rendant
        inversible un petit ens. de flèches, soit en \struck{\ill{}}
        \add{égalisant} \uncertain{des couples de flèches}, soit en
        ajoutant un petit ens. de flèches
      \item[(3)] $R_{\alpha} = \varinjlim_{\beta < \alpha} R_{\beta}$ si
        $\alpha$ ordinal limite.)
    \end{enumerate}
    \marginal{en travers de la marge, en diagonale : \ill{}}
\end{enumerate}

\underline{NB 1}\ \struck{\ill{}} \note{une ligne biffée} se pose
seulement si $(d', d'')$ n'est pas « petit », ou du moins ne \uncertain{se
ramène} à un cas petit $\mathcal{E}$\ldots\ Néanmoins, même si $(d', d'')$
petits, la question de caractériser les $I \to R$, à partir d'une
construction transfinie, se pose. Dans la définition récurrente des
théories algébriques \add{(cas d)}, il faut admettre des $\varinjlim$
sur des ens. bien ordonnés \ldots\ Il faudrait néanmoins prouver que
rien à des $I \to (R_{\lambda})_{\lambda \in \Lambda}$ satisfaisant la
condition envisagée (disons « $I$ $\Delta$-engendre $R_{\lambda}$ »)
dans l'un \ill{} pourrait pas \ill{} \uncertain{amalgamés} des
$R_{\lambda}$ sous $I$ dans $\Delta$.

\underline{NB 2}\ Soit $\tilde{I}$ la (petite) sous-catégorie pleine de
$R$ engendrée par les objets $\in I$ et les \add{(image dans $R$ des)}
sources et buts de toutes les flèches rajoutées au cours de la
construction transfinie. Il est plausible que $\forall C \in
\mathrm{Ob}\,\Delta$,
\[
  \mathrm{Hom}_{\Delta}(R, C) \longrightarrow
  \mathrm{Hom}_{\mathrm{Cat}}(\tilde{I}, C) = T_{\tilde{I}}(C)
\]
\note{au-dessus de la flèche : $\simeq T_{R}(C)$, pour le membre de
gauche} est \underline{pleinement fidèle}. Or $T_{\tilde{I}}$ est une
théorie algébrique. Il est plausible \add{dès lors} que \struck{\ill{}}
$T_{R}$ se déduit de $T_{\tilde{I}}$ \struck{\ill{}} en rajoutant
seulement des axiomes. Mais il ne s'ensuit peut-être pas de la forme en une

\page{63}
seule fois, mais qu'il faille remonter une construction transfinie.

Il y a donc des énoncés possibles pour caractériser les $R \in
\mathrm{Ob}\,\Delta$ tels que $T_{R}$ soit une théorie
\underline{algébrique} :
\begin{enumerate}
  \item[1)] $\exists$ \underline{petite} catégorie $J$, et ensembles
    $\Sigma$, \struck{$\Sigma'$} de flèches de $\tilde{J}$, tels que $R$
    se déduit de $\tilde{J}$ en rendant inv. les flèches de $\Sigma$
    (et égalisant \add{couples de} flèches de $\Sigma'$.)
    \marginal{en regard, encerclé : $+$ ens. $\Sigma'$ de
    \uncertain{couples de} flèches de $\tilde{J}$}
  \item[2)] \struck{$\exists$ petite catégorie $J$, et \ill{}} $\exists$
    \struck{famille} \add{pseudo-syst. inductif} transfini
    $(R_{\alpha})$ dans $\Delta$, tels que
    \begin{enumerate}
      \item[(i)] $R_{0} = \tilde{J}$, $J$ petite cat.
      \item[(ii)] $R_{\alpha+1}$ se déduit de $R_{\alpha}$ en rendant
        inversibles un \uncertain{petit} ens. de flèches, \underline{et}
        égalisant un ens. \uncertain{petit} de couples de flèches
      \item[(iii)] $R_{\alpha} = \varinjlim_{\beta < \alpha} R_{\beta}$
        si $\alpha$ ordinal limite
      \item[(iv)] $R_{\alpha} = R$ pour $\alpha$ assez grand, $\alpha \in
        \mathcal{U}$ (l'Univers).
    \end{enumerate}
\end{enumerate}

\begin{enumerate}
  \item[(2)] \underline{Autre Pb} : conditionner les théories
    algébriques par des propriétés intrinsèques (plutôt que par
    « constructibilité » via une suite transfinie de constructions, ou
    la représentabilité par un $R$ tel que \ldots). La seule condition
    (probablement pas suffisante) que je vois est : $T$ commute aux
    2-limites projectives de $\Delta$, i.e. : $T$ transforme
    \struck{\ill{}} pl. fid (fid.) en \uncertain{idem},
    $\mathrm{Hom}_{\mathrm{Cat}}(J, C) \to \mathrm{Hom}_{\mathrm{Cat}}(J,
    T(C))$ \ldots
    \note{le membre de gauche est écrit sans $T$ ; le contexte demande
    $T(\mathrm{Hom}_{\mathrm{Cat}}(J, C))$}
  \item[(3)] \underline{Conditions \struck{\ill{}} d'exactitude
    généralisées} (comme suggérées par Schanuel) : on donne
    \struck{\ill{}} $\rho \to \sigma$ dans $R$ et on demande sur $C$ que
    $\mathrm{Hom}_{\Delta}(\sigma, C) \to \mathrm{Hom}_{\Delta}(\rho, C)$
    soit une équivalence (resp. \add{pl.} fidèle, resp. fidèle ?\ldots).
    Cela inclut p. ex. rendre des flèches égales (cf. \uncertain{l'exemple
    1-2}). Les remarques sur l'existence de $\varprojlim$ et
    $\varinjlim$ dans $\Delta$ ainsi obtenues restent applicables.
\end{enumerate}

\page{64}
\begin{enumerate}
  \item[(4)] Changement de $\Delta \subset \Delta'$
\end{enumerate}

\begin{tikzcd}[column sep=small, row sep=small]
  \text{Thé}_{\Delta'} \arrow[rr, "\text{restriction}"] \arrow[d] & & \text{Thé}_{\Delta} \arrow[d] \\
  \mathrm{Thal}_{\Delta'} \arrow[rr, "\text{restr.}"] \arrow[d, no head, "\wr\|" description] & & \mathrm{Thal}_{\Delta} \arrow[d, no head, "\wr\|" description] \\
  \Delta'_{\mathrm{al}} \arrow[rr] & & \Delta_{\mathrm{al}}
\end{tikzcd}

2-foncteur \uncertain{d'où $\Delta'$-enveloppe}. C'est défini ainsi
\[
  \Delta' \to \Delta \qquad \rho_{\Delta', \Delta}
\]

\note{à droite du tableau, le cube dont il est la face de devant ; les
verticales doubles portent une flèche dans chaque sens}

\begin{tikzcd}[column sep=small, row sep=small, nodes={font=\scriptsize}]
  \mathrm{Th}_{\Delta'} \arrow[rrr] \arrow[dr] \arrow[dd, bend right=20] & & & \mathrm{Th}_{\Delta} \\
  & \mathrm{Thal}_{\Delta'} \arrow[rrr] & & & \mathrm{Thal}_{\Delta} \arrow[ul] \arrow[dd, bend left=20] \\
  \Delta' \arrow[rrr] \arrow[uu, bend right=20] & & & \Delta \\
  & \Delta'_{\mathrm{al}} \arrow[ul] \arrow[uu] \arrow[rrr] & & & \Delta_{\mathrm{al}} \arrow[ul] \arrow[uu, bend left=20]
\end{tikzcd}

\page{65}
\underline{Existence de $\varprojlim$ et $\varinjlim$ dans $T(C)$}
\add{$\simeq \mathrm{Hom}_{\Delta}(R, C)$}, pour $T$ 2-représentable.

Soit $\delta$ un type de diagrammes, tel que
\begin{enumerate}
  \item[a)] Les $\varprojlim$ (finies) de type $\delta$ existent dans
    $C$
  \item[b)] Les $\varprojlim$ de type $\delta$ commutent, \struck{pour tt
    $C' \in \mathrm{Ob}\,\Delta$} \add{dans $C$}, aux $\varinjlim$ de
    type $d$ (i.e. aux $\varinjlim$ de type $d''$)
\end{enumerate}
Alors les $\varprojlim$ de type $\delta$ existent dans $T(C)$, et se
calculent dans $\mathrm{Hom}_{\Delta}(R, C)$ « argument par argument »
(en d'autres termes, les foncteurs $r_{c} : T(C) \to C$ associés aux
\uncertain{constructions $\Gamma_{c}$} commutent aux $\varprojlim$ de type
$\delta$).

De plus, pour un hom de théories $T \to T'$, $T(C) \to T'(C)$ commute aux
$\varprojlim$ de type $\delta$. Idem pour $C \to C'$, $C'$ satisfaisant
les mêmes hyp. que $C$ : $T(C) \to T(C')$ commute aux $\varprojlim$ de
type $\delta$.

Donc, pour $R$, \struck{\ill{}} $\in \Delta$, \struck{\ill{}} et $J \in
\mathrm{Ob}\,\mathrm{Cat}$, $\mathrm{Hom}_{\mathrm{Cat}}(J, R) \in \Delta$
(au moins si $J$ petite)
\[
  \mathrm{Hom}_{\Delta}(R, \mathrm{Hom}_{\mathrm{Cat}}(J, C)) \simeq
  \mathrm{Hom}_{\mathrm{Cat}}(J, \mathrm{Hom}_{\Delta}(R, C))
\]
i.e.
\[
  \boxed{\ T_{R}(\mathrm{Hom}_{\mathrm{Cat}}(J, C)) \simeq
  \mathrm{Hom}_{\mathrm{Cat}}(J, T_{R}(C))\ }
\]
\note{les lettres $J$ et $R$ de ces deux lignes sont repassées et
surchargées ; la lecture $J$ est celle du contexte}

Faisons p. ex. $C = (\mathrm{Ens})$ (supposant $(\mathrm{Ens}) \in
\mathrm{Ob}\,\Delta$ !), $J$ remplacé par $J^{\circ}$ :
\[
  T_{R}(\hat{J}) \simeq \mathrm{Hom}(J^{\circ}, T_{R}(\mathrm{Ens}))
\]

Si $C \to C'$ dans $\Delta$ est pl. fidèle (resp. fid.) alors $T(C) \to
T(C')$ l'est aussi (pour $T$ 2-repr.). \struck{et} Si $T \to T'$ est
associé à $R' \to R$ « \add{$\Delta$-}générateur » (tel que, donc, la
sous-catégorie pleine $\in \Delta$ de $R'$ par laquelle se factorise
soit \struck{\ill{}} $\cong R$] alors le diagramme

\begin{tikzcd}
  T(C) \arrow[r, hook] \arrow[d] & T(C') \arrow[d] \\
  T'(C) \arrow[r, hook] & T'(C')
\end{tikzcd}

est 2-cartésien.

\page{66}
[Cas particulier où $T' = \tau^{I}$ \ldots
\note{le feuillet ne porte que ce début de crochet}

\page{67}
\begin{enumerate}
  \item[(4)] \note{le numéro, encerclé, est surchargé} $d'' = \emptyset$ :
    structures alg. définies seulement par $\varprojlim$ : Faits
    particuliers :
    \begin{enumerate}
      \item[1°)] $\exists$ $\varprojlim$ quelc. dans les $T(C)$, les
        $T(C) \to T(C')$ y commutent, de même les $T(C) \to T'(C)$
        ($T$, $T'$ 2-repr.), ainsi que les $T(C) \to C$ (foncteurs
        constructions).
      \item[2°)] $C \to \hat{C}$ est un $\Delta$-foncteur,
        \struck{\ill{}} (si $(\mathrm{Ens}) \in \mathrm{Ob}\,\Delta$)
        donc on trouve que (si $T$ 2-repr. par $R$)
        \[
          T_{R}(C) \longrightarrow T_{R}(\hat{C}) \simeq
          \mathrm{Hom}_{\mathrm{Cat}}(C^{\circ}, T_{R}(\mathrm{Ens}))
        \]
        \note{sous $T_{R}(\mathrm{Ens})$ : $S$, avec un signe de
        répétition ; les indices $R$ sont surchargés}
        est pl. fidèle, et identifie $T_{R}(C)$ à la sous-catégorie str.
        pleine de $\mathrm{Hom}_{\mathrm{Cat}}(C^{\circ}, S)$ formée des
        $\varphi : C^{\circ} \to S$ tels que, pour \ill{} \ill{}
        \underline{construction} $r \in R$, \add{d'où $F : S \to
        (\mathrm{Ens})$}, $F \circ \varphi : C^{\circ} \to (\mathrm{Ens})$
        soit repr. Et il suffit d'exiger cette condition pour \ill{}
        \ill{} $r_{i}$ qui $\Delta$-engendrent (i.e. on a un diagramme
        2-cartésien

        \begin{tikzcd}[column sep=small]
          T(C) \arrow[r] \arrow[d] & \mathrm{Hom}_{\mathrm{Cat}}(C^{\circ}, T(\mathrm{Ens})) \arrow[d, "\varphi \mapsto (\tilde{r}_{i} \circ \varphi)"] \\
          C^{I} \arrow[r] & \mathrm{Hom}_{\mathrm{Cat}}(C^{\circ}, (\mathrm{Ens}))^{I}
        \end{tikzcd}

    \end{enumerate}
\end{enumerate}

\page{68}
Donc $T$ est connu (à équiv. près) quand on connaît $S$ et les $r_{i} :
S \to (\mathrm{Ens})$ [et on connaît alors $T$ sur $\tau^{I}$\ldots], ou
de façon plus intrinsèque, si on connaît \add{l'image de $R$ dans}
\struck{\ill{}} $\mathrm{Hom}(S, \mathrm{Ens}) = \widehat{S^{\circ}}$.

\begin{enumerate}
  \item[3°)] On a $S \cong \mathrm{Hom}_{\Delta}(R, (\mathrm{Ens}))
    \subset \widehat{R^{\circ}}$, et cette sous-catégorie
    \underline{contient} $R^{\circ}$ [car les foncteurs repr. \add{cov.}
    commutent aux $\varprojlim$ !) donc
    \[
      \boxed{\ R^{\circ} \xrightarrow{\ \partial\ } S\ }
    \]
    \note{la flèche est un crochet d'inclusion, marqué $\partial$ et
    $\xi$ au-dessus}
    et le foncteur \add{$\tilde{r} : S \to (\mathrm{Ens})$,} associé à
    un $r \in \mathrm{Ob}\,R$ est le foncteur que le représente
    \underline{dans} $S$. Ainsi $R$ est connu quand on connaît $S$ et
    \uncertain{son image} ens. des $\mathrm{Hom}(S, \mathrm{Ens}) =
    \widehat{S^{\circ}}$, car $R \to \struck{\mathrm{Hom}}\,
    \widehat{S^{\circ}}$ est le composé
    \[
      R \xrightarrow{\ \partial^{\circ}\ } S^{\circ} \longrightarrow
      \widehat{S^{\circ}}
    \]
    de deux foncteurs pl. fid.

    Ainsi $S = T(\mathrm{Ens})$ est muni d'une sous-catégorie str. pleine
    $\Sigma$, ($\approx R^{\circ}$) [formée des $\varphi \in
    \mathrm{Ob}\,S$ tels que pour $\forall C \in \mathrm{Ob}\,\Delta$ et
    tout $\xi \in T(C)$, définissant $\tilde{\xi} : C^{\circ} \to S$,
    les composés $\mathrm{Hom}(\varphi, \tilde{\xi}(x))$ \add{$x
    \mapsto$} de $C^{\circ}$ dans Ens soit représentables] \add{la
    sous-catégorie} $\Sigma$ est stable dans $S$ par $\varinjlim$
    (\underline{inductives}) de type $d'$, \struck{et pour tout les
    \ill{} ainsi} \ill{} un équivalence \add{explicite} entre
    $\Sigma^{\circ}$ et $R$.
    \marginal{d'où un foncteur $T(C) \to C$, \ill{}}
  \item[4°)] Soit \struck{\ill{}} $\Theta$ l'une de toutes les
    catégories où les \ill{} petits diagrammes, $\Delta_{0}$ l'une
    \ill{} catégorie \ill{}
    \note{ce qui suit est enfermé dans une accolade et traversé d'une
    longue diagonale}
    que dans la définition de $\Delta$ ne \ill{} les $\varprojlim$ de
    type $\Theta'$ existent, Supposons ne figurent que des conditions de
    finitude, et \uncertain{sorte} que $\Delta_{0} \subset \Delta$
    \add{les}. Donc la théorie $T$ \add{$= T_{\Theta R}$} repr. pour
    \struck{\ill{}} pour $\Delta_{0}$, dont nous allons montrer
    \struck{$C \in \mathrm{Ob}\,\Delta_{0}$, \ill{}} Si c'est
    représentable, on \struck{$\mathrm{Hom}_{\Delta_{0}}(R, C) \simeq
    \mathrm{Hom}_{\varinjlim}(R^{\circ}, C^{\circ})^{\circ} \simeq$}
    soit que la catégorie $R_{0}$ qui représente dans une catégorie avec
    $\varinjlim$ \struck{les} quelconques \struck{\ill{}} et
    \note{le feuillet s'arrête ici}
\end{enumerate}

\section{Exemples}
\note{titre souligné en tête de la p. 69, de sa main}

\page{69}
\begin{enumerate}
  \item[(1)] $d' = d'' = \emptyset$. \struck{Alors} Prenant $\Delta =
    (\mathrm{Cat}_{d})$, les Pbs 1) et 2) ont une solution évidente
    resp. bien connue. Les théories algébriques sont exactement celles
    qui sont représentables par des \underline{petites} catégories. Les
    $I \to R$ générateurs sont \ill{} tous ens. surjectifs.
  \item[(2)] \underline{Conditions d'exactitudes}. On peut par exemple
    prendre $\Delta = (\uncertain{gr.-groupoïdes})$. Le Pb 1 a une
    solution \ill{} (\uncertain{construction} comme groupoïde fondamental
    d'une catégorie), le Pb 2) est trivial. Les théories alg. sont celles
    repr. par des petites \struck{catégories} \add{gr.-groupoïdes}. Les
    $I \to R$ générateurs sont les ens. surj.
    \marginal{$\rho = [a \xrightarrow{\ u\ } b]$, on rend $u$ inv.}
  \item[(3)] Prenons $\Delta = (\text{catégories ordonnées})$. Pb 1 a
    une solution évidente (catégorie ordonnée associée à une catégorie),
    Pb 2 une solution facile. Les théories algébriques sont celles
    repr. par petites cat. ordonnées. Les $I \to R$ générateurs sont
    ens. surjectifs.
    \marginal{$\rho = [a \rightrightarrows b]$, flèches $u$, $v$ ; on
    rend $u$ égal à $v$}
\end{enumerate}

\page{70}
\begin{enumerate}
  \item[(2)] $d''$ $=$ tous les « petits » diagrammes, \add{$d'' =
    \emptyset$.} Pas de conditions d'exactitude.
    \note{ce (2) reprend la numérotation des « Compléments » plutôt que
    celle des exemples ; le (1) de la p. 69 lui répond}

    \underline{Pb 1}\ $I$ une petite catégorie, on a pour $C \in
    \mathrm{Ob}\,\Delta$
    \[
      \mathrm{Hom}(I, C) \simeq \mathrm{Hom}(I^{\circ}, C^{\circ})^{\circ}
      \simeq \mathrm{Hom}_{\Delta^{\circ}}(\widehat{I^{\circ}},
      C^{\circ})^{\circ} \simeq \mathrm{Hom}_{\Delta}
      (\widehat{I^{\circ}}{}^{\circ}, C)
    \]
    \note{sous le second membre : « sans $\varinjlim$ quelc. » ; l'indice
    $\Delta^{\circ}$ du troisième est surchargé} donc $I$ a la
    $\Delta$-enveloppe $\tilde{I} = \mathrm{Hom}(I, C)^{\circ}$ (qui est
    l'opposé d'un topos). \note{le $C$ de $\mathrm{Hom}(I, C)^{\circ}$
    est ce qui est écrit ; la chaîne précédente demande
    $\widehat{I^{\circ}}{}^{\circ}$} Notons que $I \to \tilde{I}$ définit
    les objets de $I$ comme famille \underline{cogénératrice} (donc
    $I^{\circ} \to \tilde{I}^{\circ}$ comme famille génératrice).

    \underline{Pb 2}\ $C \in \mathrm{Ob}\,\Delta$, $\Sigma \subset
    \mathrm{Fl}\,C$ [$\Sigma$ petit]. Rendre inversibles les flèches de
    $\Sigma$ ? \add{est défini, montrons que c'est une solution (au sens
    de Cat)}. \struck{Supposons que} On peut supposer que $\Sigma$
    contient les isom, et tel que \marginal{$\Sigma$ un ens. plus petit
    après solution !} $\Sigma$ est stable par composition, pour chaque
    \ill{} $vu \in \Sigma$, $v \in \Sigma \Rightarrow u \in \Sigma$. Mais
    alors $\Sigma$ donne lieu à un « calcul de fractions à g. », et il
    est connu \add{dans $C\Sigma^{-1}$} que les $\varprojlim$ \add{des
    types des petites $\varprojlim$} qui existent dans $C$ existent dans
    $C\Sigma^{-1}$, et $C \to C\Sigma^{-1}$ y commute. Il reste à examiner
    les produits (infinis). Pour ceci, on peut supposer aussi que
    $\Sigma$ est stable par produits \underline{infinis} de morphismes.
    Mais alors on voit tout de suite que les produits de $C$ sont des
    produits de $C\Sigma^{-1}$. OK !
    \marginal{Mais, il faut \struck{\ill{}} \uncertain{voir} que la
    catégorie obtenue est une $\mathcal{U}$-catégorie}

    Notons que si $I$ est une famille « cogénératrice par mono strict »
    dans \struck{\ill{}} $C$ [pour tt $X \in \mathrm{Ob}\,C$, $X
    \xrightarrow{\ \sim\ } \varprojlim_{X/C} X'$], alors il en est
    encore des \uncertain{$\varprojlim$} dans $C\Sigma^{-1}$.
    \struck{Donc on voit}

    Supposons qu'on ait un \add{pseudo-}syst. inductif filtrant
    $(R_{\alpha})$ dans $\Delta$, avec $I \to R_{\alpha}$ cogénérateur.
    Soit $R = \varinjlim_{\alpha}^{\Delta} R_{\alpha}$. Alors $I \to R$
    est-il cogénérateur ? Soit $R' = \varinjlim_{\alpha}^{\mathrm{Cat}}
    R_{\alpha}$. \struck{Donc $R'$, les $\varprojlim$ \ill{} des
    \ill{}} \note{une ligne et demie biffées} Donc $R$ s'obtient à partir
    de $R'$ en rendant inversibles certaines flèches, dont $\mathrm{Ob}\,
    R \xrightarrow{\ \sim\ } \bigcup_{\alpha} \mathrm{Im}\, R_{\alpha}$.
    De ceci on déduit aussitôt ce qu'on a dit.
    \marginal{Mais, attention : le fait d'être strictement cogénérateur
    n'est pas transitif}

    Il faut aussi étudier les cas où on « rajoute des flèches » à $R$
    muni de $I$, d'où $R \to R'$. Ce $\Delta$-foncteur \struck{\ill{}}
    est ens. surjectif, et \ill{} ce qu'on veut !
\end{enumerate}

\page{71}
\underline{Th.}\ Les théories algébriques pour $\Delta$ sont celles qui
sont 2-représentables par des $R \in \mathrm{Ob}\,\Delta$ admettant des
petites familles \add{str.} \underline{cogénératrices}. Pour que
$I \xrightarrow{\ \rho\ } R$ \struck{\ill{}} \note{une ligne biffée} soit
\add{$\Delta$-}basique, il f. et s. que ses images soient \add{str.
$\Delta$-}cogénératrices (c'est plus faible que \add{str.}
cogénératrices !).
\marginal{en haut à droite, relié au mot « cogénératrices » : passe bien
dans la construction transfinie, \ill{} \ill{} par \ill{}}

La nécessité dans la première assertion est claire par ce qui précède.
Pour la suffisance, \struck{on peut supposer} \uncertain{résultat} des
résultats de la deuxième assertion. Or soit $I'$ la catégorie \add{muni
d'un foncteur pl. fid. $I' \to R$} (telle que $\mathrm{Ob}\,I' = I$,
$\mathrm{Ob}\,\varphi = \varphi_{0}$. \struck{Considérons} On sait que
$T_{\tilde{I}'}$ est une théorie \underline{algébrique} de base $I$. Or
$\widehat{I'^{\circ}} \to R^{\circ}$ \add{et surjectif sur les Hom}
puisque $\rho : I' \to R$ est ens. surjectif, donc $\tilde{I}' \to R$
aussi. Donc il \struck{suffit de} \struck{\ill{}} \ill{} à rendre
\uncertain{iso} des \struck{\ill{}} certaines flèches des \uncertain{iso}.
\note{la fin de la phrase est surchargée ; le sens est que $R$ se déduit
de $\tilde{I}'$ en rendant inversibles certaines flèches}

\underline{Remarques}\ \begin{enumerate}
  \item[(1)] On \struck{procède} \add{obtient} ici à la construction les
    étapes successives : 1°) Objets de base $I$ 2°) rajouter des flèches
    3°) axiomes d'égalité sur des \add{couples de} flèches (on a
    maintenant $T_{\tilde{I}'}$), 4°) \struck{axiomes} \add{rendre
    inversibles des flèches, ou encore} \struck{d'égalité sur d'autres
    flèches, ou encore} \note{ces derniers mots portent des surcharges qui
    ne se laissent pas démêler} On peut d'ailleurs condenser 3° et 4°) en
    une seule étape, \ill{} (lorsque des axiomes successives ne se font
    \ill{} en une seule fois, grâce au calcul de fractions \add{comme
    catégorie de fractions, \ill{}} \uncertain{topos}) par un ens. de
    flèches satisfaisant les conditions de restriction déjà dites.
  \item[(2)] Soit $R \in \mathrm{Ob}\,\Delta$. Les catégories quotients
    $R' \in \mathrm{Ob}\,\Delta$ déduites de $R$ par rendre inversibles
    des flèches \struck{sont exactement celles qui} \add{correspondent
    1-1 aux} ens. de flèches $\Sigma$ \add{$\Sigma$} telles que
    \begin{enumerate}
      \item[a)] $\Sigma$ stable par composition, par changement de base,
        contient les iso
      \item[b)] $\Sigma$ stable par produits quelconques.
    \end{enumerate}
  \item[(3)] Supposons que $R$ ait une petite famille
    \underline{cogénératrice} $I$. Alors les $\mathrm{Hom}_{\Delta}(R,
    \mathrm{Ens}) \simeq R^{\circ}$ [les foncteurs de $R$
\end{enumerate}

\page{72}
dans (Ens) qui commutent aux $\varprojlim$ sont les foncteurs
représentables), i.e. si $T = T_{R}$, on a $R = S^{\circ}$, où $S =
T(\mathrm{Ens})$. \struck{\ill{}} \struck{Donc, \ill{}} ou encore
$R^{\circ} = S$. Donc :

\underline{Cor}\ Les \add{2-}théories algébriques définissables par
$\varprojlim$ \add{(sans axiomes d'exactitude)} \struck{correspondent aux
catégories} sont 2-équiv. par $T \mapsto T(\mathrm{Ens})$ à celle des
$\mathcal{U}$-catégories $S$ avec $\varprojlim$ quelc. et
\struck{familles} \add{petites} sous-ens. génératrices par épi. stricts
[i.e. telles que $S \to \widehat{I^{\circ}}$ soit conservatif]. La donnée
d'une base \add{de $T$} équivaut à celle dans \ill{} d'une famille
d'objets cogénératrice de $S$.
\note{le corollaire est marqué d'un trait vertical en marge}

Pour toute $C \in \mathrm{Ob}\,\Delta$, on a
\[
  T(C) \simeq \mathrm{Hom}_{\Delta}(S^{\circ}, C)
\]
L'équivalence est obtenue en notant que le deuxième est la catégorie
formée des foncteurs tels que $\forall X \in C$, le composé avec
$\mathrm{Hom}_{C}(X, -)$ est représentable, donc le deuxième s'interprète
comme les foncteurs $S^{\circ} \times C^{\circ} \to (\mathrm{Ens})$ qui
sont représentables par chaque argument fixé, ou comme les foncteurs
$C^{\circ} \to S$ qui sont tels que pour tout $Y \in S$, sa composé avec
$\mathrm{Hom}(Y, -)$ soit représentable. Or on a
\[
  T(C) \longrightarrow \mathrm{Hom}(C^{\circ}, T(\mathrm{Ens})) =
  \mathrm{Hom}(C^{\circ}, S)
\]
(déduit de $C^{\circ} \to \mathrm{Hom}_{\Delta}(C, \mathrm{Ens}) \subset
\widehat{C^{\circ}}$ induisant $C^{\circ} \to \mathrm{Hom}(T(C),
T(\mathrm{Ens}))$, i.e. $T(C) \to \struck{\ill{}}
\mathrm{Hom}(C^{\circ}, T(\mathrm{Ens}))$).

\begin{enumerate}
  \item[(4)] Comme $S = T(\mathrm{Ens})$ a des $\varprojlim$ quelc. et
    $S \cong R^{\circ}$, $R$ a aussi des $\varinjlim$ quelc., donc
    $R \to R'$ \add{$\varprojlim$ et $\varinjlim$} entre objets de
    $\Delta$. \note{la ligne est écrite en surcharge sur une amorce
    biffée}
  \item[(5)] Soit un foncteur $\varphi : R \to R'$ \add{$S$ et $R$ ont
    toutes les $\varprojlim$}. Si $R$ a une petite famille
    cogénératrice, alors $\varphi$ est un $\Delta$-foncteur, i.e.
    commute aux $\varprojlim$, ssi il admet un adjoint à gauche $\psi$
    \[
      \mathrm{Hom}(\psi(Y'), X) \simeq \mathrm{Hom}(Y', \varphi(X))
    \]
    Alors $\varphi$ est un \uncertain{plongement} à une sous-catégorie
    \struck{$R'$ est une sous-catégorie $\psi$ est} \ill{} \ill{} ssi
    $\psi$ est pl. fidèle, (et alors $R'$ a une p. famille cog.).
    \struck{Ainsi les sous-catégories} \struck{\ill{}} Les théories de
    la th. algébrique $R$ \add{$R_{0}$} correspondent \struck{aux sous}
    à : \underline{certaines} sous-catégories pleines de $R$, savoir
    celles qui ont des $\varprojlim$ et pleines par des foncteurs
    d'inclusion $\varphi$ ont un adjoint à gauche. En termes de $S =
    R^{\circ} = T(\mathrm{Ens})$, \add{et $S' = R'^{\circ} =
    T'(\mathrm{Ens})$} \struck{ce sont les sous-catégo-}
    \marginal{ainsi les $\Delta$-foncteurs $\varphi : R \to R'$ (i.e.
    $T' \to T$) correspondent aux foncteurs $\psi : S' \to S$ \ill{} qui
    \uncertain{admettent} un adjoint \ill{} ; \ill{} \ill{}}
\end{enumerate}

\page{73}
\note{en tête : $S' = R'^{\circ}$, $S = R^{\circ}$, au-dessus des deux
membres de la ligne suivante}
ries $\varphi^{\circ} : T'(\mathrm{Ens}) \to T(\mathrm{Ens})$ est
simplement le foncteur associé à $T' \to T$ déduit de $\varphi$. Ce
foncteur doit donc commuter aux $\varprojlim$ quelc. [donc ce qui
exprime simplement que $\psi$ commute aux $\varinjlim$, \struck{alors}
et c'est un adjoint à g.] Donc les \uncertain{sous-}théories de $T$
correspondent aux sous-catégories pleines $S'$ de $S$, \struck{admettant
des} [stables par $\varprojlim$, en un certain sens \ill{} par les
autres conditions] admettant des \underline{$\varprojlim$} quelconques
(mais pas nécessairement stables par les dites) et telles que le
foncteur d'inclusion $S' \xrightarrow{\ i\ (=\varphi^{\circ})\ } S$ ait
un adjoint à \underline{gauche} $j$ ($= \psi^{\circ}$)
\[
  \mathrm{Hom}_{S}(x, i(y')) \simeq \mathrm{Hom}(jx, y')
\]
\struck{Alors on a des foncteurs \ill{} adjoints \ill{}} \struck{[Mais en
général ? \ill{} (Structure (foncteurs $S'$ $S$-structure)}
\note{deux lignes biffées, une amorce de crochet abandonnée} $S' \to$
structure « libre » définie par une $S$-structure) Mais en général, il
n'y a pas moyen de définir pour tt $C \in \mathrm{Ob}\,\Delta$ un
foncteur $T(C) \to T'(C)$ qui les redonne \uncertain{pour} $C =
\mathrm{Ens}$ [\underline{NB}\ Cependant, si $C$ est un
\underline{topos}, et si $T$, $T'$ se définissent au \uncertain{st.}
des $\varprojlim$ finies, c'est OK \ldots]

\underline{NB}\ La 2-catégorie des théories algébriques pour $\Delta$
est 2-équivalente à la 2-catégorie des cat. $S$ qui sont stables par
$\varprojlim$ \add{petites} \struck{quelc.} et qui ont une petite
sous-catégorie str. génératrice, dans les foncteurs qui commutent aux
$\varprojlim$ (mais, en renversant les flèches pour les morphismes ?).
\struck{Ce sont aussi les catégories telles que les foncteurs qui
transforment $\varprojlim$ quelc. en \ill{} \ill{}} \note{une ligne et
demie biffées}

\page{74}
\begin{enumerate}
  \item[(3)] $d'$ \struck{toutes les} \add{$d'$} \struck{famille de
    diagrammes} \add{quelconque}, \struck{\ill{}} \struck{les \ill{}}
    finis $d'' = \emptyset$, pas de \underline{conditions
    d'exactitude}. $d'_{0}$ l'ens. de tous les diagrammes finis,
    $d'_{1}$ l'ens. de tous les petits diagrammes, et $\Delta \supset
    \Delta_{0} \supset \Delta_{1}$ les solutions correspondantes.
    \note{les indices $0$ et $1$ des deux $d'$ se lisent mal ; l'ordre
    $\Delta \supset \Delta_{0} \supset \Delta_{1}$ demande que $d'_{0}$
    soit le plus petit}

    \underline{Pb 1}\ $\Delta$-Enveloppe \add{$\rho_{\Delta}(I)$} d'une
    catégorie $I$. Soit $\bar{I}$ \struck{\ill{}} \add{$= \widehat{I^{\circ}}{}^{\circ}$}
    sa $\Delta_{1}$-enveloppe, et soit $\tilde{I}$ la
    \struck{sous-}catégorie pl. de $\bar{I}$ formée des
    $\varprojlim$ des types $d'$ d'objets de $I$ [\underline{NB}\ $I \to
    \bar{I}$ pl. fid. !]. Je dis que cela marche, \underline{pourvu que}
    \ill{} \uncertain{les objets} \ill{} $\tilde{I}$ est stable pour
    les $\varprojlim$ des diagrammes finis discrets, ou $d' \ni$ un $d'$
    \ill{} [$C$ c'est le cas si $d' =$ ens. des \ill{} de tous les
    diagrammes finis (!)\ldots) sinon on doit prévoir et prendre pour
    $\tilde{I}$ la plus petite sous-catégorie \add{str.} pleine de
    $\bar{I}$ contenant $I$ et stable par $\varprojlim$ de type $d'$.
    \marginal{en diagonale, dans la marge : pb 1 et pb 1' \ill{} \ill{}
    \uncertain{petits} \ill{} \uncertain{diagrammes finis}}
    En effet, on a pour $C \in \mathrm{Ob}\,\Delta$
    \[
      \mathrm{Hom}_{\struck{\mathrm{Cat}}}(I, C) \longrightarrow
      \mathrm{Hom}_{\Delta_{1}}(\bar{I}, \bar{C}),
    \]
    \note{le premier Hom porte un indice biffé et repassé « Cat » ; le
    second est marqué $\Delta_{1}$} et l'iso. \uncertain{gén.} \add{de
    foncteurs [pl. fidèles \ill{}]} \uncertain{fournie} \add{de(s)
    foncteurs $\bar{I} \xrightarrow{\ \bar{\varphi}\ } \bar{C}$ tels que
    $\bar{\varphi}(I) \subset \struck{\ill{}} C$ (ess.)}. Mais
    \struck{\ill{}} Soit $\rho(C)$ la $\Delta$-$\Delta_{1}$ enveloppe de
    $C$ i.e. la $\Delta_{1}$-catégorie \add{$\varphi'(C)$} avec
    $\Delta$-foncteur $C \xrightarrow{\ i\ } \rho(C)$
    \underline{universel} [c'est une catégorie pleine de \uncertain{foncteurs}
    de $\bar{C}$]. \struck{Admettons que} \struck{le foncteur est
    pleinement fidèle (Pb auxiliaire)} \struck{$I \to C$ se prolonge
    en foncteurs} \note{une ligne et demie biffées} Alors les foncteurs
    $\varphi : I \to C$ \struck{En effet, \ill{} les foncteurs
    \ill{} \ill{} à $\varphi$ les f. \ill{} $I$ dans l'image
    ess.} \uncertain{correspondent aux} foncteurs $\psi = i\varphi :
    \tilde{I} \to \rho(C)$ \add{qui \uncertain{envoient} $I$ dans
    l'image essentielle de $C$ par $i$}, ou encore les foncteurs
    $\bar{\psi} : \bar{I} \to \rho(C)$ tels que $\bar{\psi}(I)$ soit
    contenu dans l'image essentielle de $C$ ; de tels $\bar{\psi}$
    envoient $\tilde{I}$ dans $C$, d'où
    \[
      \mathrm{Hom}_{\mathrm{Cat}}(I, C) \longrightarrow
      \mathrm{Hom}_{\Delta}(\tilde{I}, C),
    \]
    donc un foncteur « restriction » en sens inverse. Il est clair que
    les composés dans les deux sens sont l'identité, OK.

    \underline{Pb 1'}\ \underline{ou} Pb \underline{auxiliaire}\ Si $C
    \in \mathrm{Ob}\,\Delta$, montrer que $C \to \rho(C)$ est pl.
    fidèle. Je dis qu'on a
    \[
      \rho(C) = \mathrm{Hom}_{\Delta}(C, \mathrm{Ens})^{\circ} =
      T_{C}(\mathrm{Ens})^{\circ},
    \]
\end{enumerate}

\page{75}
le foncteur $C \to \rho(C)$ \add{$= \mathrm{Hom}(C, \mathrm{Ens})^{\circ}
= \widehat{C^{\circ}}{}^{\circ}$} \struck{donné} \uncertain{donné} par
\uncertain{un} associant : tt objet de $C$ le foncteur covariant qu'il
représente (donc il est clair que $C \to \rho(C)$ est pl. fidèle !). On
a en effet, si $D$ est une \struck{catégorie où les $\varprojlim$ quelc.
\ill{} et si $\rho(C)$ est définie comme le} \add{$\Delta \to
\Delta_{1}$ enveloppe}
\begin{gather*}
  \mathrm{Hom}_{\Delta_{1}}(C, D) \simeq \mathrm{Hom}_{\Delta_{1}}
  (\rho(C), D) \simeq \ill{}, \\
  \rho(C) \cong T_{\rho(C), \Delta_{1}}(\mathrm{Ens})^{\circ} \cong
  T_{C, \Delta}(\mathrm{Ens})^{\circ} = \mathrm{Hom}_{\Delta}(C,
  \mathrm{Ens})^{\circ} : \text{OK.}
\end{gather*}
\note{les deux lignes sont reliées par des traits d'insertion ; l'ordre
donné ici est celui que les traits indiquent}

\underline{Pb 2}\ $C \in \Delta$, $\Sigma \subset \mathrm{Fl}$,
construire $C\Sigma^{-1}$ \underline{dans $\Delta$}, et des ens. de
couples de flèches qu'on veut égaliser. \struck{Tout marche bien si}
\struck{quotient} \add{d'où} Soit $\rho(C, \Sigma)$ la
\add{$\Delta_{1}$-}catégorie de \add{fractions} de $\rho(C)$, définie
par $\Sigma$, et soit $C'$ la plus petite \add{$\Delta$-}sous-catégorie
\add{str. pleine} de $\rho(C, \Sigma)$ \struck{qui} telle que $C \to
\rho(C, \Sigma)$ se factorise par $C'$, \add{i.e. telle que $C'$
contienne la sous-catégorie \underline{pleine} $C\Sigma^{-1}$}. (Soit
$D \in \mathrm{Ob}\,\Delta$, alors les $\Delta$-foncteurs $\varphi : C
\to D$ correspondent aux $\Delta_{1}$-foncteurs $\psi : \rho(C) \to
\rho(D)$ tels que $\psi$ envoie $C$ dans (la ss-cat. pleine) $D$, et
les $\Delta$-$\Sigma$-foncteurs « i.e. » tels que les foncteurs $\psi' :
\rho(C, \Sigma) \to \rho(D)$ \struck{\ill{}} qui envoient $C$ dans $D$,
\struck{ou \ill{}} ou qui envoient \uncertain{aussi} $C'$ dans $D$. Donc
\[
  \mathrm{Hom}_{\Delta, \Sigma}(C, D) \xrightarrow{\ \alpha\ }
  \mathrm{Hom}_{\Delta}(C', D),
\]
\note{sous la flèche, une flèche en sens inverse marquée $\beta$} et
\uncertain{$\beta$} le foncteur restriction en sens inverse. Il est clair
que les deux composés \ill{} $\beta\alpha$ \struck{\ill{}} les foncteurs
identiques, \ill{} \uncertain{pour} $\alpha\beta$ on note que
\struck{\ill{} $\beta$ est pl. fidèle} $\beta$ est pl. fidèle. OK. C'est
déduit de $C\Sigma^{-1} \subset \rho(C, \Sigma)$ en rajoutant
transfiniment des $\varprojlim$ de type $d'$, et que sait que
$\mathrm{Hom}_{\Delta}(C', D) \to \mathrm{Hom}_{\mathrm{Cat}}(C, D)$ est
pl. fidèle (car $\mathrm{Hom}_{\Delta}(C', D) \to
\mathrm{Hom}_{\mathrm{Cat}}(C', D)$ l'est \ldots).
\marginal{en diagonale, dans la marge, sur deux registres : i.e. $\Sigma$
\uncertain{contient} \ill{} les diagrammes \ill{} \uncertain{produits}
et \ill{} \ill{} dans $C\Sigma^{-1}$, les \ill{} \ill{} ; \ill{} p. ex.
\ill{} finis discrets}

\underline{NB}\ Dans les deux cas envisagés (Pb 1 et Pb 2), la
catégorie résultante (de $I$, resp. de $C$) se déduit de la
sous-catégorie pleine de $I$ (resp. de $C$) \ill{} en rajoutant des
$\varprojlim$ de type $d'$ (transfiniment au besoin). On trouve sauf
erreur :
\marginal{en regard, en diagonale : il faudrait encore regarder \ill{}
des $\varprojlim$ \ill{} \ill{} \uncertain{fibres}}

\underline{Th.}\ Les théories algébriques de type $\Delta$ sont celles
qui sont 2-représentées par des $R \in \mathrm{Ob}\,\Delta$ ayant une
petite sous-catégorie \add{$I'$} \struck{qui} $d'$-génératrice [i.e.
\struck{tous objets de} $R$ est la réunion d'une suite transfinie de
sous-catégories str. pleines de $R$ déduites en rajoutant des
$\varprojlim$ de type $d'$]. Un foncteur $I \to R$ \add{($I$ discret)}
est basique ssi il est $d'$-cogénérateur.

\page{76}
On est encore ramené à prouver ce dernier énoncé. On procède comme dans
2°), en introduisant la sous-cat. pleine $I'$ de $R$ telle que
$\mathrm{Ob}\,I' = I$, d'où $T_{R} \to T_{\tilde{I}'}$ \underline{pl.
fidèle}. Regardons $\tilde{I}' \to R$, il est \add{ess.} \struck{surjectif
sur les objets et surjectif sur} les Hom, \struck{\ill{}} est-il ens.
surjectif ? OK. Si \struck{la fonction} $d'$ \struck{\ill{}} satisfait
aux conditions écrites \uncertain{plus haut}, car alors on trouve un
\add{$\Delta$-}foncteur pl. fidèle
\[
  \tilde{I}'\Sigma^{-1} \longrightarrow R,
\]
\note{sous $\tilde{I}'\Sigma^{-1}$ : « catégorie de fractions
\ill{} »} \struck{\ill{}} qui \uncertain{donc} est ens. surjectif. Donc
on trouve

\begin{tikzcd}[column sep=small, row sep=small]
  & \tilde{I}' \arrow[dr] \arrow[ddl] & \\
  I' \arrow[ur] \arrow[rr] \arrow[d, "\text{pl. fid.}"'] & & \bar{I}' \arrow[d, "\text{ess. surj., et surjectif sur les Hom}"] \\
  R & & \rho(R) \arrow[ll, hook']
\end{tikzcd}

$I'$ cogénérateur strict $R$ \\
qui cogénérateur str. $\rho(R)$ \\
donc $I'$ cogénérateur str. $\rho(R)$ (?)
\note{ces trois lignes sont disposées en escalier sous le diagramme ;
le point d'interrogation est le sien}

Si $d'$ satisfait aux conditions envisagées (p. ex. $d' =$ ens. des
diagrammes finis)

\underline{Cor\,\uncertain{Main}}\ On peut construire les théories
algébriques pour $\Delta$ en trois pas, comme pour $\Delta_{1}$ : on
prend une catégorie $I'$ ayant $I$ comme ens. d'objets (pas 1° et 2°) et
prend une catégorie des fractions de $\tilde{I}'$.

\note{marqué d'un double trait en marge :} Mais la conclusion est
valable si p. ex. $d'$ est formé des diagrammes discrets, stable par
passage à un sous-diagramme ( \struck{\ill{}}

Néanmoins, dans le cas d'un $d'$ général ($=$ formé de diagrammes
finis) il semble qu'il faille une construction récurrente transfinie

\underline{Question}\ \struck{\ill{}} Supposons $d'$ petit. Alors une
$R$ \struck{telle que $R$} admettant une petite ss-catégorie
$d'$-cogénératrice (donc \struck{2-rep} correspondant à une théorie
algébrique pour $\Delta$) est petite, car \ldots
\note{les membres de la phrase sont reliés par des traits de renvoi ;
l'ordre donné est celui des traits}

\page{77}
\begin{enumerate}
  \item[4°)] Cas $d' =$ ens. des diagrammes finis, $d'' = \emptyset$,
    pas de \underline{conditions d'exactitude}.

    Comme $d'$ est petite, les $R \in \mathrm{Ob}\,\Delta_{0}$ qui
    2-représentent des th. algébriques sont les $R$ ens. petites
    (stables par $\varprojlim$ finies). \struck{\ill{}} \note{une ligne
    et demie biffées}
    \begin{quote}
      $T$ théorie algébrique pour $\Delta$ \\
      $R \in \mathrm{Ob}\,\Delta$ qui représente $T$ \\
      $S = T(\mathrm{Ens}) = \underline{\mathrm{Hom}}_{\mathrm{lex}}(R,
      \mathrm{Ens})$ \\
      $\Sigma \subset S$ (structures représentables)
    \end{quote}
    \note{les quatre lignes sont réunies par une accolade}
    Tout se détermine mutuellement
    \begin{align*}
      R :\quad & T = T_{R} : C \mapsto
        \underline{\mathrm{Hom}}_{\mathrm{lex}}(R, C) \\
      & S = \struck{\ill{}}\ \underline{\mathrm{Hom}}_{\mathrm{lex}}(R,
        \mathrm{Ens}) \\
      & \Sigma \cong R^{\circ} \subset S \text{ par foncteurs repr.} \\
      \Sigma :\quad & R = \Sigma^{\circ} \\
      & T : C \mapsto \underline{\mathrm{Hom}}_{\mathrm{lex}}
        (\Sigma^{\circ}, C) \\
      & S = \underline{\mathrm{Hom}}_{\mathrm{lex}}(S^{\circ}, C),
        \text{ contient } \Sigma \text{ comme sous-catégorie pleine.}
    \end{align*}
    \note{la dernière ligne est écrite ainsi, avec $S^{\circ}$ et $C$ ;
    la ligne symétrique demande $\Sigma^{\circ}$ et Ens ; un mot biffé
    suit}

    \struck{Reste à voir \ill{}} Autre façon de déduire $S$ de $R$ ou
    $\Sigma$

    \underline{Lemme} : \struck{\ill{}} Le foncteur inclusion $\Sigma \to
    S$ définit une \underline{équivalence}
    \[
      \mathrm{Ind}(\Sigma) \longrightarrow S
    \]
    \marginal{en diagonale, dans la marge : cf. SGA 4 I 8.7.5 (a)
    $\Sigma \in S \subset \underline{\mathrm{Hom}}(\ill{}, \mathrm{Ens})$
    \ill{} ; (ii) $\Sigma$ stable par $\varprojlim$ finies ; (iii)
    $\Sigma$ \ill{} \uncertain{petits} \ill{}}
    En effet, \struck{\ill{}} on sait que $\Sigma \to S$ est universel
    pour des foncteurs $\Sigma \to D$ \add{\ill{} $D$ à $\varinjlim$
    \uncertain{filtrantes}} exacts à gauche. Or montrons que $\Sigma \to
    \mathrm{Ind}(\Sigma)$ \add{(On sait que $\mathrm{Ind}(\Sigma)$ a des
    $\varinjlim$ \uncertain{filtrantes} et des $\varprojlim$ finies
    (donc quelc.) car $\Sigma$ les a, et $\Sigma \to \mathrm{Ind}(\Sigma)$
    y commute.)} \struck{En effet les} \marginal{SGA 4 I 8.9.5 ;
    puis, plus bas : SGA 4 I 8.7.3} \uncertain{Les} foncteurs
    \underline{exacts} $\Sigma \to D$ correspondent aux foncteurs
    $\mathrm{Ind}(\Sigma) \to D$ qui commutent aux $\varinjlim$
    \add{\uncertain{filtrantes}} et aux $\varprojlim$ finies, soit ceux
    qui, sur $\Sigma$, y commutent [nécessité triviale, pour la
    suffisance, $\mathrm{Ind}(\Sigma) \to \mathrm{Ind}(D)$ \struck{commute
    aux} est exact à droite, \struck{\ill{} finies (\ill{}, \ill{},
    \ill{})} (c'est un adjoint à g, qui commute à \ill{}) et par
    \uncertain{I 8.9.8} $\mathrm{Ind}\,D \to D$ l'est aussi (c'est
    \struck{aux limites inductives}) ok]

    \underline{Corollaire}\ \struck{Pour tout} \add{Soit $X \in S$. Pour
    que $X \in \Sigma$, il f. et s. que} \add{$X \in \Sigma$,} le
    foncteur $\mathrm{Hom}(X, -)$ sur $S$ commute aux $\varinjlim$
    filtrantes.
    \note{la phrase est récrite par-dessus une première rédaction ; les
    membres sont donnés dans l'ordre des traits de renvoi}

    C'est évidemment nécessaire, par construction des Hom dans
    $\mathrm{Ind}(S)$. C'est suffisant, car \struck{si $X \in S$,} on
    aura $X = \varinjlim_{i} X_{i}$ ($X_{i} \in \mathrm{Ob}\,\Sigma$),
    \add{filtrantes}, \uncertain{le morphisme} identique se factorise
    par un des $X_{i}$, donc \ill{} un $X_{i} \to X$ ayant une section,
    donc $X$ est un \uncertain{composant} \struck{\ill{}} \add{p. ex.}
    \ill{} donc stable (c'est un noyau) \struck{donc $\Sigma \to S$}
    OK.
    \marginal{en bas à gauche, en diagonale : NB Montrer \ill{}
    \uncertain{catégorie} \ill{} ? cf. aussi SGA 4 I 8.7.8}
\end{enumerate}

\page{78}
Donc $\Sigma$ se reconstitue à partir de $S$ comme \struck{l'ensemble
des ob.} sous-catégorie (str.) pleine des objets « de présentation
finie » !
\begin{align*}
  S :\quad & \Sigma = S_{\mathrm{pf}} \subset S \\
  & R = \Sigma^{\circ} \\
  & T : C \mapsto \underline{\mathrm{Hom}}_{\mathrm{lex}}
    (S_{\mathrm{pf}}{}^{\circ}, C) \\
  T :\quad & S = T(\mathrm{Ens}) \\
  & \Sigma = S_{\mathrm{pf}} = T(\mathrm{Ens})_{\mathrm{pf}} \\
  & R = \Sigma^{\circ}
\end{align*}

\underline{Corollaire}\ Pour qu'une \struck{théorie} $\Delta_{1}$-théorie
\add{$T$} (par $\varprojlim$ quelc.) provienne d'une théorie
définissable par $\varprojlim$ finies, il f. et s. que $S = S_{T} =
T(\mathrm{Ens})$ soit \struck{\ill{}} strictement engendrée par ses
\add{la sous-catégorie pleine des} objets de \uncertain{prés.} finie,
ou encore que $S \to \struck{\ill{}}\ \hat{\Sigma}$ soit pleinement
fidèle. (Alors $\Sigma^{\circ}$ est la catégorie classifiante pour la
théorie sur $\Delta_{0}$ qui donne naissance à $T$).

\underline{Corollaire}\ Le 2-foncteur de restriction : $\Delta_{1}$ des
théories algébriques sur $\Delta_{0}$ est \textbf{1}-fidèle, i.e. pour
\struck{\ill{}} deux théories $T$, $T'$ sur $\Delta_{0}$,
\[
  \underline{\mathrm{Hom}}_{\mathrm{Th}_{\Delta_{0}}}(T, T')
  \longrightarrow \underline{\mathrm{Hom}}_{\mathrm{Th}_{\Delta_{1}}}
  (T_{1}, T'_{1}) \quad \text{pl. fidèle.}
\]
L'image essentielle est formée des foncteurs $T_{1} \to T'_{1}$ qui sont
tels que
\[
  S = T(\mathrm{Ens}) \add{= T_{1}(\mathrm{Ens})} \longrightarrow S' =
  T'(\mathrm{Ens}) = T'_{1}(\mathrm{Ens})
\]
transforme objets de P.F. en objets de P.F. [On pourrait les appeler les
morphismes de prés. finie, ou cohérents, de théories\ldots].

\underline{Question}\ Considérons $\Delta$ défini par $d' =$ ens. des
diagrammes finis \underline{discrets}, $d'' = \emptyset$, pas de
condition d'exactitude. Soit $R$ une petite cat. $\in \mathrm{Ob}\,
\Delta$,
\[
  \Sigma = R^{\circ} \subset S = T_{R}(\mathrm{Ens}) =
  \underline{\mathrm{Hom}}_{\mathrm{prod}}(R, \mathrm{Ens}).
\]
Peut-il se faire que $S$ soit défini par des sous-catégories
\add{str. pleines} \struck{essentielles} \ill{} \struck{\ill{} on voit
aussi} Comment trouver \add{des} $\Sigma_{0}$ telles que $S \cong
\rho_{\Delta, \Delta_{1}}(\Sigma)$ ? On aura évidemment $\Sigma \subset
S_{\mathrm{pf}} = \Sigma_{0}$ et $S_{\mathrm{pf}} \cong \rho_{\Delta,
\Delta_{0}}(\Sigma)$ (conditions nécessaires) \add{et suffisantes ?}
Donc la question se ramène en termes d'une $\Sigma_{0} \in \mathrm{Ob}\,
\Delta_{0}$ et d'une \add{sous-catégorie pleine} \struck{famille
génératrice dans $\Sigma_{0}$ de} $\Sigma$, stable par sommes
finies\ldots
\note{le paragraphe est marqué d'un trait vertical en marge ; le
« question » de tête est surchargé}

\page{80}
\marginal{en diagonale, en haut à gauche, au crayon : \uncertain{traiter}
le cas avec quelques $d'$ \uncertain{distincts}}
\note{liste de cas, numérotés dans des cercles ; des flèches en marge
renvoient (4') au-dessus de (4'') et (5') au-dessus de (5), et le
numéro de la ligne « Topos » est repassé}
\begin{enumerate}
  \item[(1)] $\varprojlim$ finies, foncteurs exacts
  \item[(1')] $\varprojlim$ de type $d'$, foncteurs qui y commutent
  \item[(1'')] Catégories \add{pré}ordonnées avec Inf finis, foncteurs
    qui y commutent
  \item[(2)] $\varinjlim$ finies et $\varprojlim$ finies, foncteurs
    exacts
  \item[(2')] Prétopos (axiomes « finis » des topos) foncteurs exacts
  \item[(2'')] Catégories semi-additives, foncteurs additifs
  \item[(2''')] Catégories additives, foncteurs additifs
  \item[(2'''')] Catégories abéliennes, foncteurs exacts [foncteurs
    additifs ?] \note{un numéro biffé précède}
  \item[(3)] $(\mathrm{Cat}_{d})$, avec $(d', d'')$ petits
  \item[(4)] Cat avec $\varprojlim$ finies, $\varinjlim$ filtrantes
    exactes, foncteurs qui y commutent \add{et $+$ \ill{} ?}
  \item[(4'')] Cat. préordonnées avec inf finis, sup quelconques, que
    inf distributifs par sup quelc., foncteurs commutant aux inf finis
    et sup quelc.
  \item[(4')] $\varprojlim$ finies, $\varinjlim$ quelc., sommes disjointes
    universelles, $\varinjlim$ filtrantes exactes, foncteurs
    \add{exacts} commutant aux $\varinjlim$
  \item[(5)] $\varprojlim$ et $\varinjlim$ quelc., foncteurs qui y
    commutent
  \item[(5')] Topos
  \item[(6')] Topos essentiels ($\cong \hat{C}$, i.e. ayant assez de pts
    essentiels, i.e. ayant les propriétés \add{d'exactitude} (pour
    $\varinjlim$ et $\varprojlim$ de la cat. des ens.), morphismes
    essentiels de topos.
\end{enumerate}
\note{le numéro de « Topos » se lit (5') ou (6) ; la dernière ligne
parle de (6) et (6')}

Je m'attends à des ennuis (de sortie de l'Univers) \add{des ennuis avec
(6) et (6')}

\end{document}
